\documentclass[12pt,a4paper]{article}
% Preâmbulo institucional comum — DEE/IFMA
\usepackage[utf8]{inputenc}
\usepackage[T1]{fontenc}
\usepackage[brazil]{babel}
\usepackage{lmodern}
\usepackage{microtype}
\usepackage{geometry}
\usepackage{xcolor}
\usepackage{graphicx}
\usepackage{booktabs}
\usepackage{tabularx}
\usepackage{array}
\usepackage{enumitem}
\usepackage{hyperref}
\usepackage{lastpage}
\usepackage{fancyhdr}
\usepackage{setspace}
\usepackage{titlesec}
\usepackage{tcolorbox}
\usepackage{ragged2e}

\geometry{a4paper,top=2.2cm,bottom=2.2cm,left=2.5cm,right=2.5cm}
\definecolor{azulpetroleo}{HTML}{0F4C5C}
\definecolor{ambar}{HTML}{C8892D}
\definecolor{cinzaclaro}{HTML}{F2F5F7}
\definecolor{cinzaescuro}{HTML}{374151}
\hypersetup{colorlinks=true,linkcolor=azulpetroleo,urlcolor=azulpetroleo,citecolor=azulpetroleo}
\setlength{\parindent}{1.25cm}
\setlength{\parskip}{0.35em}
\onehalfspacing
\titleformat{\section}{\large\bfseries\color{azulpetroleo}}{}{0pt}{}
\titleformat{\subsection}{\normalsize\bfseries\color{azulpetroleo}}{}{0pt}{}
\pagestyle{fancy}
\fancyhf{}
\fancyhead[L]{\small Departamento de Eletroeletrônica — DEE}
\fancyhead[R]{\small IFMA — Campus São Luís/Monte Castelo}
\fancyfoot[C]{\small Página \thepage\ de \pageref{LastPage}}
\renewcommand{\headrulewidth}{0.4pt}
\renewcommand{\footrulewidth}{0.4pt}
\newcommand{\campo}[1]{\textcolor{ambar}{\textbf{[#1]}}}
\newcommand{\instituicao}{Instituto Federal de Educação, Ciência e Tecnologia do Maranhão}
\newcommand{\campus}{Campus São Luís/Monte Castelo}
\newcommand{\departamento}{Departamento de Eletroeletrônica — DEE}
\newcommand{\cidade}{São Luís — MA}
\newcommand{\cabecalhoInstitucional}{%
\begin{center}
{\Large\bfseries\color{azulpetroleo}\instituicao}\\[2mm]
{\large\bfseries\campus}\\[1mm]
{\normalsize\departamento}
\end{center}
\vspace{2mm}
\hrule
\vspace{4mm}
}
\newtcolorbox{destaque}{colback=cinzaclaro,colframe=azulpetroleo,boxrule=0.6pt,arc=2mm}

\usepackage{longtable}
\begin{document}
\cabecalhoInstitucional

\begin{center}
{\LARGE\bfseries\color{azulpetroleo} Dossiê Institucional do Departamento de Eletroeletrônica}\\[2mm]
{\large Programa de Modernização Tecnológica dos Laboratórios do DEE}\\[1mm]
{\small IFMA -- Campus São Luís/Monte Castelo}
\end{center}

\tableofcontents
\newpage

\section{Apresentação}
O Departamento de Eletroeletrônica -- DEE do IFMA -- Campus São Luís/Monte Castelo atua na formação técnica e superior nas áreas de Eletrotécnica, Eletrônica, Automação e Engenharia Elétrica. Suas atividades envolvem ensino, pesquisa, extensão, iniciação científica, trabalhos de conclusão de curso, orientação acadêmica e capacitação técnica.

A modernização dos laboratórios busca aproximar ainda mais a formação acadêmica das tecnologias utilizadas pela indústria, pelo setor elétrico, por integradores de automação, empresas de telecomunicações e centros de pesquisa. A atualização dessa infraestrutura amplia as possibilidades de aulas práticas, pesquisa aplicada e desenvolvimento de projetos em parceria com o setor produtivo.

\section{Áreas de atuação}
O trabalho do Departamento está organizado em quatro grandes áreas:
\begin{enumerate}[leftmargin=*]
\item \textbf{Automação}: CLPs, IHMs, controle, instrumentação, redes industriais, sensores, robótica, IoT e Indústria 4.0;
\item \textbf{Eletrônica e Telecomunicações}: eletrônica analógica e digital, sistemas embarcados, instrumentação eletrônica, RF, redes, telecomunicações e prototipagem;
\item \textbf{Eletrotécnica}: instalações elétricas, medidas, máquinas, comandos, acionamentos, manutenção e qualidade da energia;
\item \textbf{Sistemas de Energia}: geração, transmissão, distribuição, proteção, subestações, smart grid, geração distribuída e análise de sistemas elétricos.
\end{enumerate}

% Corpo docente e distribuição por macroárea do DEE

\section{Corpo docente}\label{sec:corpo}

O DEE reúne \textbf{34 professores}, com atuação distribuída entre as principais macroáreas do Departamento. Essa composição permite integrar ensino, pesquisa aplicada, extensão e desenvolvimento tecnológico.

\begin{center}
\begin{tcolorbox}[
  colback=white,
  colframe=azulpetroleo!35,
  boxrule=0.6pt,
  arc=1.5mm,
  left=3mm,right=3mm,top=3mm,bottom=3mm,
  width=\textwidth
]
\begin{minipage}[t]{0.28\textwidth}
\centering
{\sffamily\fontsize{28}{30}\selectfont\bfseries\color{azulpetroleo} 34}\\[1mm]
{\sffamily\large\bfseries docentes}\\[1mm]
{\small\color{cinzamedio} composição total do Departamento}
\end{minipage}
\hfill
\begin{minipage}[t]{0.66\textwidth}
\small
\textbf{Atuações docentes por macroárea}\\[-0.2mm]
{\footnotesize\color{cinzamedio}Respostas com sobreposição entre áreas}\\[1.5mm]

\noindent\textbf{Automação} \hfill \textbf{10}
\begin{tikzpicture}[baseline=(current bounding box.center)]
\fill[azulpetroleo!15] (0,0) rectangle (10,0.30);
\fill[azulpetroleo] (0,0) rectangle (7.1,0.30);
\end{tikzpicture}\\[1.5mm]

\noindent\textbf{Eletrônica e Telecomunicações} \hfill \textbf{14}
\begin{tikzpicture}[baseline=(current bounding box.center)]
\fill[azulpetroleo!15] (0,0) rectangle (10,0.30);
\fill[azulpetroleo] (0,0) rectangle (10,0.30);
\end{tikzpicture}\\[1.5mm]

\noindent\textbf{Eletrotécnica} \hfill \textbf{14}
\begin{tikzpicture}[baseline=(current bounding box.center)]
\fill[azulpetroleo!15] (0,0) rectangle (10,0.30);
\fill[azulpetroleo] (0,0) rectangle (10,0.30);
\end{tikzpicture}\\[1.5mm]

\noindent\textbf{Sistemas de Energia} \hfill \textbf{6}
\begin{tikzpicture}[baseline=(current bounding box.center)]
\fill[azulpetroleo!15] (0,0) rectangle (10,0.30);
\fill[azulpetroleo] (0,0) rectangle (4.3,0.30);
\end{tikzpicture}
\end{minipage}
\end{tcolorbox}
\end{center}

{\small\color{cinzamedio}Alguns docentes atuam em mais de uma área; por isso, os quantitativos por macroárea não são exclusivos e sua soma é maior que o total de 34 professores.}

A composição multidisciplinar do corpo docente permite integrar competências tradicionais da eletroeletrônica a tecnologias digitais emergentes em atividades de ensino, pesquisa, extensão e P\&D.


\section{Formação profissional, pesquisa e P\&D}
Os laboratórios atendem os cursos técnicos de Eletrotécnica, Eletrônica e Automação e o curso de Engenharia Elétrica. Também são utilizados em iniciação científica, TCCs, monografias e pesquisas vinculadas a mestrado e doutorado, inclusive em cooperação com outras instituições.

Com equipamentos atuais, os estudantes podem trabalhar com tecnologias próximas das encontradas no ambiente profissional, desenvolvendo experiência em configuração, operação, ensaio, medição, diagnóstico, supervisão e manutenção de sistemas reais. Isso contribui para uma formação mais sólida e prepara profissionais capazes de atuar com segurança e domínio técnico.

O DEE também está disponível para desenvolver, em conjunto com empresas e instituições parceiras, projetos de \textbf{Pesquisa e Desenvolvimento (P\&D)}, pesquisa aplicada, provas de conceito, protótipos, ensaios, medições e validações tecnológicas em áreas relacionadas às competências do Departamento.

\section{Estrutura laboratorial}
A modernização contempla diferentes ambientes e tecnologias, permitindo que o apoio de cada parceiro seja direcionado às áreas em que sua atuação tenha maior aderência.

\begin{longtable}{p{0.31\textwidth} p{0.61\textwidth}}
\toprule
\textbf{Eixo/Laboratório} & \textbf{Tecnologias prioritárias}\\
\midrule
\endhead
Automação Industrial & CLPs, IHMs, I/O remoto, SCADA, redes industriais, fontes e sensores\\
Eletrônica e Sistemas Digitais & Osciloscópios, fontes, geradores, DMMs, LCR, analisadores lógicos e kits embarcados\\
Telecomunicações/RF & Analisadores de espectro, VNA, RF, comunicações móveis, fibra óptica e redes\\
Máquinas e Acionamentos & Motores, transformadores, inversores, soft-starters, proteção e bancadas de ensaio\\
Instalações e Medidas & Proteção BT, instrumentos de medição, megômetros, analisadores e segurança elétrica\\
Qualidade de Energia & Analisadores de potência, registradores, harmônicos e monitoramento\\
Proteção e Sistemas de Energia & Relés, IEDs, RTAC, IEC 61850, TC/TP, testes e automação de subestações\\
Instrumentação Industrial & Pressão, vazão, nível, temperatura, 4--20 mA, válvulas e calibração\\
Eletrônica de Potência & MOSFET, IGBT, SiC, GaN, drivers, conversores e controle digital\\
Pneumática/Hidráulica & Bancadas didáticas, válvulas, atuadores, sensores e controle\\
Computação para Engenharia & Workstations, servidores, monitores, software e aquisição de dados\\
\bottomrule
\end{longtable}

\section{Integração dos laboratórios}
% Fluxo em blocos da estrutura laboratorial do DEE.
% Layout 2x2 com eixo central: o distribuidor "MACROÁREAS" desce por um corredor
% livre entre as colunas e alimenta as quatro caixas por derivações curtas,
% sem contornar o desenho e sem sobreposição entre caixas.

\begin{center}
\resizebox{0.96\textwidth}{!}{%
\begin{tikzpicture}[
  >=stealth,
  topo/.style={
    draw=azulpetroleo,
    fill=azulpetroleo!9,
    very thick,
    rounded corners=2mm,
    align=center,
    text width=10.4cm,
    minimum height=1.20cm,
    inner sep=3mm,
    font=\sffamily\bfseries
  },
  centro/.style={
    draw=azulpetroleo!70,
    fill=white,
    thick,
    rounded corners=1.8mm,
    align=center,
    text width=3.5cm,
    minimum height=0.85cm,
    inner sep=2.5mm,
    font=\sffamily\bfseries\small
  },
  area/.style={
    draw=azulpetroleo!80,
    fill=cinzaclaro,
    thick,
    rounded corners=2mm,
    align=left,
    text width=7.0cm,
    minimum height=4.0cm,
    inner sep=4mm,
    font=\small
  },
  uso/.style={
    draw=ifmaverde!75!black,
    fill=ifmaverde!6,
    thick,
    rounded corners=2mm,
    align=center,
    text width=15.8cm,
    minimum height=1.35cm,
    inner sep=3.5mm,
    font=\footnotesize
  },
  eixo/.style={thick,color=azulpetroleo!75},
  seta/.style={->,thick,color=azulpetroleo!75},
  no/.style={circle,fill=azulpetroleo!75,inner sep=0pt,minimum size=1.6mm}
]

% Bloco principal
\node[topo] (dee) at (0,0) {Departamento de Eletroeletrônica -- DEE\\[-0.2mm]
\normalfont\footnotesize IFMA -- Campus São Luís/Monte Castelo};

% Nó de distribuição
\node[centro] (macro) at (0,-1.8) {MACROÁREAS};
\draw[seta] (dee.south) -- (macro.north);

% Primeira linha (centro em y = -5,0; corredor central livre de 1,4 cm)
\node[area] (auto) at (-4.6,-5.0) {
\centering\sffamily\bfseries AUTOMAÇÃO\par\vspace{2mm}
\raggedright\normalfont
$\bullet$ Automação Industrial\\[0.8mm]
$\bullet$ Redes Industriais e IoT\\[0.8mm]
$\bullet$ Robótica\\[0.8mm]
$\bullet$ Pneumática e Eletropneumática\\[0.8mm]
$\bullet$ Instrumentação e Sensoriamento};

\node[area] (elet) at (4.6,-5.0) {
\centering\sffamily\bfseries ELETRÔNICA E TELECOMUNICAÇÕES\par\vspace{2mm}
\raggedright\normalfont
$\bullet$ Eletrônica Analógica e Digital\\[0.8mm]
$\bullet$ Sistemas Digitais e Embarcados\\[0.8mm]
$\bullet$ Telecomunicações, RF e Redes\\[0.8mm]
$\bullet$ Instrumentação Eletrônica\\[0.8mm]
$\bullet$ Eletrônica de Potência};

% Segunda linha (centro em y = -10,2; 1,2 cm de folga para a primeira linha)
\node[area] (eletrot) at (-4.6,-10.2) {
\centering\sffamily\bfseries ELETROTÉCNICA\par\vspace{2mm}
\raggedright\normalfont
$\bullet$ Instalações Elétricas\\[0.8mm]
$\bullet$ Medidas Elétricas\\[0.8mm]
$\bullet$ Máquinas Elétricas\\[0.8mm]
$\bullet$ Comandos Elétricos\\[0.8mm]
$\bullet$ Acionamentos e Qualidade de Energia};

\node[area] (sep) at (4.6,-10.2) {
\centering\sffamily\bfseries SISTEMAS DE ENERGIA\par\vspace{2mm}
\raggedright\normalfont
$\bullet$ Proteção e Subestações\\[0.8mm]
$\bullet$ Sistemas Elétricos de Potência\\[0.8mm]
$\bullet$ Smart Grid\\[0.8mm]
$\bullet$ Geração Distribuída\\[0.8mm]
$\bullet$ Análise e Monitoramento de Sistemas};

% Uso compartilhado
\node[uso] (uso) at (0,-14.0) {\textbf{Uso compartilhado da infraestrutura:} ensino técnico e superior $\bullet$ aulas práticas $\bullet$ iniciação científica $\bullet$ TCCs e monografias $\bullet$ pesquisas de mestrado e doutorado $\bullet$ extensão $\bullet$ capacitação profissional $\bullet$ projetos de P\&D, prototipagem, ensaios e validação tecnológica.};

% Eixo central: desce pelo corredor entre as colunas e termina no uso compartilhado.
\draw[eixo] (macro.south) -- (0,-12.9);
\draw[->,thick,color=ifmaverde!65!black] (0,-12.9) -- (uso.north);

% Derivações curtas para cada macroárea.
\foreach \y/\esq/\dir in {-5.0/auto/elet,-10.2/eletrot/sep}{%
  \node[no] at (0,\y) {};
  \draw[seta] (0,\y) -- (\esq.east);
  \draw[seta] (0,\y) -- (\dir.west);
}

\end{tikzpicture}%
}
\end{center}


Os ambientes são utilizados de forma compartilhada. Um mesmo equipamento pode atender diferentes disciplinas, cursos e projetos, além de apoiar pesquisas e atividades de desenvolvimento tecnológico.

\section{Possibilidades de parceria}
O apoio pode ocorrer por meio de doação de equipamentos, kits didáticos, unidades de demonstração, softwares, licenças, treinamento, modernização de bancadas ou cooperação técnica. Também são de interesse equipamentos descontinuados comercialmente, desde que estejam em boas condições de funcionamento e possam ser utilizados com segurança em atividades acadêmicas.

A parceria pode ainda envolver workshops, capacitações, estudos de caso, TCCs, dissertações, teses e projetos de P\&D relacionados às tecnologias da empresa. O DEE pode acompanhar a utilização dos recursos recebidos e apresentar resultados acadêmicos e técnicos decorrentes da cooperação.

\section{Resultados esperados}
A modernização dos laboratórios deve ampliar a qualidade das atividades práticas, fortalecer a pesquisa aplicada e aumentar o contato dos estudantes com tecnologias utilizadas no setor produtivo. Espera-se também ampliar a capacidade do Departamento para desenvolver projetos com empresas, apoiar pesquisas de pós-graduação e contribuir para a formação de profissionais mais preparados para atuar nas áreas de eletroeletrônica.

Toda parceria ou doação será formalizada pelos setores competentes do IFMA, de acordo com os procedimentos institucionais aplicáveis.

\section{Contato}
\textbf{Francisco dos Santos Viana}\\
Chefe do Departamento de Eletroeletrônica -- DEE\\
IFMA -- Campus São Luís/Monte Castelo\\
\href{mailto:francisco.viana@ifma.edu.br}{francisco.viana@ifma.edu.br}\\
(98) 98816-6263

\vfill
\begin{center}
São Luís -- MA, 2026.
\end{center}
\end{document}
